\frametitle{The Hierarchy, and One Two-Level Step}
\footnotesize

Uniform refinement of the torus mesh gives a nested sequence of finite element spaces
$V_0\subset V_1\subset\dots\subset V_L$, with $\dim V_\ell=n_\ell$ and
$h_\ell=h_0 2^{-\ell}$. Between consecutive levels we use the canonical prolongation and
the transposed restriction, with the Galerkin coarse operator:
\[
P_\ell=P_{\ell\leftarrow\ell-1}:\R^{n_{\ell-1}}\to\R^{n_\ell},
\qquad
R_\ell:=P_\ell^{\mathsf T},
\qquad
A_{c,\ell}=P_\ell^{\mathsf T}A_\ell P_\ell .
\]
Since $A_\ell$ is SPD and $P_\ell$ has full column rank, $A_{c,\ell}$ is SPD as well, and
the coarse problem is the restriction of the fine problem to the coarse space rather than
an independently assembled operator. For the rest of this part the level index is fixed
and dropped: $A:=A_\ell$, $P:=P_\ell$, $A_c:=P^{\mathsf T}AP$.

\vspace{1mm}
One two-level step is: \emph{relax, form the residual, restrict it, solve the coarse
problem for the correction, prolongate and add, relax again}. Writing $\widetilde u$ for
the current approximation, $r=b-A\widetilde u$ for the residual and
$e=A^{-1}r=u-\widetilde u$ for the algebraic error, the coarse correction $e_H$ solves
$A_c e_H=Rr$, so
\[
\widetilde u^{\text{new}}=\widetilde u+PA_c^{-1}R\,r ,
\qquad
e^{\text{new}}=e-PA_c^{-1}RAe
=\underbrace{\bigl(I-PA_c^{-1}RA\bigr)}_{=:\,C_h}e .
\]

\begin{block}{}
$C_h$ is what one two-level step does to the error, in one line. The next frame says it is
not merely \emph{close to} a projection: with the exact coarse solve and the transposed
restriction used here, it \emph{is} exactly the operator the error splitting is built
from --- and that is why the analysis that follows is possible at all.
\end{block}
